\documentclass[10pt]{article}
\usepackage[margin=20mm]{geometry}
\usepackage{amsmath,amssymb,fontspec,unicode-math,xurl,hyperref}
\IfFontExistsTF{DejaVu Serif}{\setmainfont{DejaVu Serif}}{\setmainfont{Latin Modern Roman}}
\setmathfont{latinmodern-math.otf}
\setlength{\parindent}{0pt}\setlength{\parskip}{6pt}\sloppy
\title{Technical Verification: One Leaf Reveals the Season: Occlusion-Based Contrastive Learning with Semantic-Aware Views for Efficient Visual Representation}
\author{AI-assisted technical review}
\date{}
\begin{document}
\maketitle
Abstract: The cross-view objective in Algorithm 1 is well defined, but the displayed loss has incompatible indexing and the accompanying concentration explanation reverses the parameter direction. Static inspection of the frozen release identifies material recipe differences and a distinction between optimizer batch size and the negative pool used by the loss. Independent calculations establish these local discrepancies. They do not determine which implementation produced the reported benchmarks, and no central empirical failure is established.\par
\section*{Paper identity and scope}
One Leaf Reveals the Season: Occlusion-Based Contrastive Learning with Semantic-Aware Views for Efficient Visual Representation. Xiaoyu Yang; Lijian Xu; Hongsheng Li; Shaoting Zhang. ICML 2025, PMLR 267, I25-43759. The reviewed source is the frozen published proceedings PDF. Page numbers below are one-based physical PDF pages.\par
The main text, Tables 1–12 and Appendices A–B were read; Figures 1–3 and source pages 2–9 and 16 were actually viewed. The selected author source is https://github.com/XiaoyuYoung/OCL, commit `fd2cbdbfa430edc3059a1f3555929bada172ac50`. Static tracing covered masking, encoding, the loss, the pre-training entry point and loop, the launcher and documented recipe, with a targeted check of linear-probe freezing. Author code was neither imported nor executed. This is a method and source-consistency review, not an ImageNet reproduction.\par
\section*{Problem 1: the displayed denominator is not the implemented objective}
Equation reference: eq:correct. Section 2.3 (p. 4) defines \(B\) image pairs, one normalized vector \(u_i\) and \(v_i\) from each image, and \(B^2-B\) cross-view negatives. Algorithm 1 forms a \(B\times B\) similarity matrix and applies cross-entropy in both directions. Its row loss is \[
 \ell_i=-\tau s(u_i,v_i)+\log\sum_{k=1}^{B}\exp\{\tau s(u_i,v_k)\},\] with the positive term included in the denominator. Equation (3), however, sums to \(2B\) while retaining a lower-branch vector \(y_k^L\), and excludes \(k=i\). There are only \(B\) such lower-branch vectors. With the pair ordering used in Algorithm 1, \(k=i\) is the positive rather than the anchor’s self-comparison.\par
These differences affect the calculation. For \(B=2\), let \(u_1=v_1=e_1\) and \(u_2=v_2=e_2\), choose the paper’s \(\kappa=64\) and an admissible positive scale \(\tau=1\). Both vectors have unit norm. The positive and negative similarities are \(1\) and \(1/130\). Equation [eq:correct] (\#eq:correct) gives \[\ell_i=\log(1+e^{-129/130})=0.3153363.\] Merely changing the upper limit to \(B\) while retaining the exclusion gives \(-129/130=-0.9923077\), a different objective; if the scale grows, this latter expression tends to \(-\infty\) rather than zero. The expression with the original upper limit is undefined under branch-local indexing.\par
A charitable global-index interpretation would concatenate both branches and exclude the anchor itself. That produces an ordinary all-view contrastive objective, but also introduces same-branch negatives absent from Algorithm 1. It requires changing the vector notation and still does not describe the supplied \(B\times B\) computation. The frozen `models\_mae.py`, lines 167–189, agrees with Algorithm 1 and includes the diagonal positives. The supported correction is Equation [eq:correct] (\#eq:correct) and its transpose average. This is a Minor formula defect: the paper itself contains the coherent intended algorithm, and no loss error in the historical experiments is established.\par
\section*{Problem 2: the concentration explanation has its direction reversed}
With cosine \(c\in[-1,1]\) and \(\kappa\ge0\), Equation (2) gives \[s_\kappa(c)=\frac{1+c}{2[1+\kappa(1-c)]}.\] The following paragraph says that decreasing \(\kappa\) makes the function more concentrated around nearby features. Direct differentiation and a half-height calculation give \[\frac{\partial s_\kappa(c)}{\partial\kappa}
 =-\frac{1-c^2}{2[1+\kappa(1-c)]^2},\qquad
 s_\kappa(c)\ge\tfrac12\ \Longleftrightarrow\ c\ge\frac{\kappa}{1+\kappa}.\] The derivative is strictly negative for \(-1<c<1\), while \(s_\kappa(1)=1\). Thus increasing \(\kappa\) narrows the region above half the peak. At \(\kappa=1/64\) its cosine threshold is \(1/65\), whereas at \(\kappa=64\) it is \(64/65\). An independent finite-difference check at \((c,\kappa)=(.3,4)\) agrees with the analytic derivative to \(2.1\times10^{-11}\).\par
Section 3.2.4 (p. 7) instead describes increasing \(\kappa\) as making the task more demanding; this is useful counterevidence to a systematic conceptual error. The mathematical correction is local: reverse the direction in the p. 4 explanation. Positive-pair labels remain assigned by image identity rather than by a similarity threshold. Neither this derivative nor the half-height criterion proves the proposed semantic interpretation or predicts the optimum in Table 6. Assessment: Minor.\par
\section*{Problem 3: the frozen release does not specify the reported model recipe}
The paper specifies a fixed Xavier-initialized patch projection (Section 2.1, p. 2), a trainable scale in the loss (p. 4), and \(\kappa=64\) by default (p. 5). In the frozen `models\_mae.py`, lines 26–68 and 162–164, the default is \(\kappa=1/64\) and the temperature is the ordinary scalar \(0.07\). Similarities use the same rational formula as Equation (2), followed by division by that scalar. The pre-training constructor in `main\_pretrain.py`, line 157, passes no concentration override, and the argument parser exposes none.\par
This is not a harmless reciprocal notation for \(\kappa\): the formula multiplies by the stored value in both sources. For orthogonal vectors, \[s_{64}(0)=1/130=0.0076923,\qquad
 s_{1/64}(0)=32/65=0.4923077.\] At \(c=.5\) the values are \(0.0227273\) and \(0.7441860\). Their relation is not a single global rescaling or additive shift that a softmax temperature could absorb. The fixed reciprocal temperature could represent the numerical value \(\tau=1/.07\), but it does not implement training of \(\tau\).\par
The projection weights are initialized and then used directly in the forward path; only the positional embedding is explicitly frozen. The pre-training optimizer receives the model’s ordinary trainable parameters, with no projection exclusion or detachment. Under the ordinary `timm.PatchEmbed` behavior, the projection therefore remains trainable. This contradicts the fixed-projection recipe; it does not assert that its resulting gradients are nonzero on every input.\par
There is a separate immediate recipe inconsistency: `PRETRAIN.md` passes `mae\_vit\_large\_patch16`, while the frozen model namespace defines the `sae\_vit\_large\_patch16` family. The entry point indexes that namespace directly, so the documented name cannot resolve if execution reaches model construction. This conclusion is static; the launcher was not run. The appendix also lists a ViT-L pre-training batch of 2,400 (Table 10, p. 16), whereas the main-text default is 2,048.\par
These are Minor release and specification defects with unresolved historical provenance. An exact benchmark configuration, including the concentration, scale parameterization, projection freeze and actual batch layout, is needed to choose a faithful repair. Simply replacing defaults here would not establish reproduction of 85.8\% accuracy. No author adoption or repaired training run is claimed.\par
\section*{Qualification 4: optimizer batch size and the contrastive negative pool}
The discussion around Table 3 (p. 6) attributes batch-size gains to a larger negative pool. `PRETRAIN.md` recommends gradient accumulation to retain an effective batch size when memory or GPU count is reduced. The entry point calculates this size as \(bAR\), where \(b\) is the per-device microbatch, \(A\) is the accumulation count and \(R\) is the number of devices.\par
The release’s loss, however, constructs similarities only from the current local input. No cross-device feature gathering appears in this path. `engine\_pretrain.py`, lines 45–60, computes the loss separately for each microbatch, divides by \(A\), and accumulates its gradients. Distributed gradient synchronization and later scalar logging do not add feature vectors to a loss denominator. Each anchor therefore has \(b-1\) negative candidates, not \(bAR-1\).\par
The distinction is exact, not merely a potential optimization difference. In the two-image example of Problem 1, processing both images together gives loss \(0.3153363\). Processing two singleton microbatches with \(A=2\) gives zero in each: the sole class has probability one. Its gradient is identically zero. For a nonstationary comparison, keep \(u_1=e_1,u_2=e_2,v_2=e_2\) and set \(v_1=(\cos\theta,\sin\theta)\) at \(\theta=.3\). The full symmetric loss is \(0.4477130\), with central-difference derivative \(0.2866228\) in \(\theta\) (step \(10^{-6}\)); the accumulated singleton derivative remains exactly zero. All vectors are normalized and both procedures have optimizer effective batch two.\par
Accumulation still increases the number of images contributing to an optimizer step, so that conventional meaning of effective batch is valid. It cannot be used to infer an equivalent contrastive pool. Nor is a local-negative objective intrinsically invalid. The documented recommendation does not explicitly promise an identical negative pool. Moreover, varying the local microbatch at fixed accumulation and device counts would enlarge both the effective batch and the local pool, a reading compatible with the Table 3 explanation. The available evidence therefore establishes a scope clarification, not a demonstrated batch-semantics error or a falsification of Table 3. An assertion that the historical runs held the pool fixed while attributing gains to a larger pool would need additional run configuration evidence. Faithful comparisons must report both microbatch/pool size and optimizer batch size; a changed global-pool implementation would require feature aggregation with appropriate gradients and its own validation.\par
\section*{Checks that survive and boundaries of the evidence}
The masking path samples a random permutation and retains an even number of tokens before splitting it. For \(196\) patches, global masking ratios \(.3\) and \(.4\) retain 136 and 116 tokens, or 68 and 58 per branch. The sets are disjoint and the visible fractions \(.34694\) and \(.29592\) agree with the rounded \(.35\) and \(.30\) descriptions. Positional embeddings are attached before shuffling. Figures 1–2 illustrate the intended pairing but do not independently establish semantic diversity for every masked image.\par
The arithmetic behind several headline summaries is consistent: the FLOP reduction from \(15.6\) to \(12.0\) is \(23.08\%\); the COCO AP-box gain over the listed MoCo-v3 baseline is \(3.6\) percentage points; the ADE20K Base-to-Large gain is \(7.1\) points. The abstract’s \(133\times4=532\) GPU-hours is compatible with a rounded 533 GPU-hours in Table 3. Figure 3 supports a reported efficiency tradeoff, not an independently standardized hardware measurement. Table 3 also shows a slight decline from 2,048 to 4,096, so a general monotonic batch-size claim would be too strong.\par
The release applies a shared random crop and flip before masking. The paper’s narrower statement about not generating separate augmented views is compatible with that path; it should not be read as absence of all preprocessing. The targeted linear-probe check freezes the backbone and trains the head. Tables 7–9 include outcomes below competing methods, so the supported empirical claim is a tradeoff rather than universal accuracy dominance. Masking and batch size co-vary in some ablations; the tables do not isolate every claimed semantic cause. No measured accuracy, runtime or robustness result is independently refuted by this review.\par
\section*{Conclusion}
Dimension | Recorded state\par
Highest supported impact | Minor\par
Scientific adjudication | Pending\par
Repair verification | Loss and concentration identities verified\par
Author adoption | Not established\par
Human reviews | None recorded\par
The highest supported impact is Minor. The displayed objective and concentration explanation need local correction, and the release needs an unambiguous configuration and batch-pool specification for faithful replication. Algorithm 1 remains coherent. The available evidence does not establish a central benchmark failure or determine the historical implementation.\par
1 Xiaoyu Yang; Lijian Xu; Hongsheng Li; Shaoting Zhang. One Leaf Reveals the Season: Occlusion-Based Contrastive Learning with Semantic-Aware Views for Efficient Visual Representation. PMLR 267 (2025), published PDF. https://proceedings.mlr.press/v267/yang25ao.html.\par
\section*{Scope and limitations}
The independently written `independent\_checks.py` evaluates the normalized-vector examples, concentration derivatives and thresholds, random-mask set properties and table arithmetic; outputs are in `independent\_check.json`. Selected source excerpts are retained in `code\_evidence.txt`. No dependency installation, author-code execution, checkpoint loading or training was performed. The source inspection is selective; full fine-tuning, detection, segmentation and robustness pipelines were not audited or reproduced. Original benchmark logs and matched historical configurations are unavailable in this packet.\par
A focused check revisited the original loss, concentration calculation, release recipe and batch-pool argument, and narrowed the last item to a scope clarification. Scientific adjudication remains pending and no human review is recorded.\par
\end{document}
